\section{Introducción}

Tecnologías como la IA generativa y las arquitecturas RAG están cambiando la manera en la que construimos aplicaciones. Esto ha revolucionado el desarrollo de software, hoy en día las aplicaciones dejaron de basarse en reglas rígidas para convertirse en sistemas complejos que de manera dinámica coordinan modelos de lenguaje, bases vectoriales y agentes inteligentes.

La presión del mercado por lanzar continuamente y rápidamente innovaciones han conducido a los desarrolladores a ignorar buenas practicas de diseño de software. Se hace frecuente encontrar código desordenado o donde encontramos mezclas de elementos que deberían de encontrarse separados (ejemplo, los prompts con lógica principal) generando así deuda técnica, como evidencia de esto, investigaciones recientes muestran que el 77\% de los errores de las aplicaciones RAG provienen de una mala arquitectura de integración y no del modelo en sí \cite{shao_are_2024}. 

Hoy en día los errores de diseño solo son percibidos una vez que el sistema falla o descubiertos de forma manual, a pesar de los riesgos que conlleva la mala estructuración. Por este motivo el presente articulo revisa la literatura actual para clasificar estos nuevos errores (antipatrones) con el fin de justificar la necesidad de crear a futuro una herramienta que los detecte automáticamente analizando el código fuente.